\section{Conditional probability, independence and Bayes' rule}
\subsection*{Core problems}
\begin{problembook}
Two standard dice are rolled independently.
\begin{enumeratea}
  \item What is the probability that the sum is less than 8, given that it is odd?
  Are the events ``the sum is less than 8'' and ``the sum is odd'' independent?
  \item Are ``the sum is 7'' and ``the first die shows 3'' independent?
  What about ``the sum is 8'' and ``the first die shows 3''?
\end{enumeratea}
\end{problembook}
\begin{problem}
Let $A$ and $B$ be events with $\Pbb(A)>0$ and $0<\Pbb(B)<1$.
\begin{enumeratea}
  \item Prove that if $A$ and $B$ are independent, then so are $A$ and $\overline{B}$,
  and also $\overline{A}$ and $\overline{B}$.
  \item Prove that the following four conditions are equivalent:
  \[
  \Pbb(A\cap B)>\Pbb(A)\Pbb(B),\quad
  \Pbb(A\mid B)>\Pbb(A),\quad
  \Pbb(B\mid A)>\Pbb(B),\quad
  \Pbb(A\mid B)>\Pbb(A\mid \overline{B}).
  \]
  \item If $A$ is independent of $B$ and $A$ is independent of $C$, must $A$ be
  independent of $B\cap C$? What about $B\cup C$?
\end{enumeratea}
\end{problem}
\begin{problembook}
Three cards are drawn successively without replacement from a deck of 52 cards.
Find the probability that:
\begin{enumeratea}
  \item all three are aces;
  \item the third card is an ace;
  \item the second card is an ace, given that the first is not.
\end{enumeratea}
\end{problembook}
\begin{problembook}
A box contains three coins: two fair coins and one with heads on both sides.
We choose a coin uniformly and toss it. Conditional on the chosen coin, successive tosses are independent.
\begin{enumeratea}
  \item The outcome is heads. What is the probability that the coin is the false one?
  \item We toss the same coin again and get heads again. Find the new probability
  in two ways: by updating the preceding answer, treating it as a new prior,
  and directly from the two tosses.
  \item What is the probability of heads on the third toss, given that the first two were heads?
\end{enumeratea}
\end{problembook}
\begin{problempencil}[Test 1, SE 2023]
An urn contains 4 yellow, 5 green and 2 blue pebbles. We select a pebble uniformly.
If it is yellow, we roll one standard die; if green, two; if blue, three.
The rolls are independent. Given that the sum is 3, what is the probability that the pebble was green?
\end{problempencil}
\begin{problembook}
A test gives a positive result for a person with a disease with probability $99\%$
(sensitivity), and a negative result for a healthy person with probability $99\%$
(specificity). The disease affects $0.5\%$ of the population. We choose a person uniformly and test them.
\begin{enumeratea}
  \item What is the probability that a person with a positive result actually has the disease?
  \item At what disease prevalence is this probability exactly $1/2$?
  \item The person is tested again and the result is again positive. If the two tests
  are conditionally independent given the person's health status, and have the same
  sensitivity and specificity, what is the probability that the person has the disease?
  When is this independence assumption questionable?
\end{enumeratea}
\end{problembook}
\subsection*{Additional problems}
\begin{probleminteresting}[Kahneman, Thinking, Fast and Slow]
An accident occurs at night in a city with a blue and a green taxi company.
$85\%$ of the taxis are green and $15\%$ are blue. Assume that a taxi's probability
of being involved in an accident does not depend on its company.
A witness states that the taxi was blue. A subsequent assessment shows that the
witness identifies the colour correctly with probability $80\%$, both for blue
and for green taxis. What is the probability that the taxi really was blue?
\end{probleminteresting}
\begin{problem}
Three tokens have sides coloured white--white, black--black and white--black,
respectively. We choose a token uniformly and place it with a randomly selected
side facing up, the two sides being equally likely. Given that the upper side is
white, what is the probability that the hidden side is also white?
\end{problem}
\begin{probleminteresting}[Monty Hall]
A car is behind one of three doors, chosen uniformly. You select a door without
knowing where the car is. The host knows its location and always opens a different
empty door, then offers you the chance to switch to the remaining closed door.
\begin{enumeratea}
  \item Should you switch? Does the probability of winning under the ``always switch''
  strategy depend on how the host chooses a door when there are two possibilities?
  \item You chose door 1 and the host opened door 3. When the host can choose between
  doors 2 and 3, they open door 3 with probability $p\in[0,1]$.
  Find the conditional probability that the car is behind door 2. For which $p$ is switching worthwhile?
\end{enumeratea}
\end{probleminteresting}
\begin{probleminteresting}[Birthday paradox]
Assume that different people's birthdays are independent and uniformly distributed
over 365 days. What is the smallest number of people for which the probability
that at least two share a birthday is greater than $1/2$? Use the multiplication
rule. Discuss the assumptions of the model.
\end{probleminteresting}
\begin{problem}
Two fair coins are tossed independently. Let $A$ be ``the first coin shows heads'',
$B$ be ``the second coin shows heads'', and $C$ be ``the coins show the same outcome''.
Prove that $A,B,C$ are pairwise independent but not mutually independent.
Are $A$ and $C$ independent conditional on $B$?
\end{problem}
\begin{probleminteresting}[Simpson's paradox]
The following table contains real data on two treatments for kidney stones
(Charig et al., \emph{BMJ} 292, 1986; Julious \& Mullee, \emph{BMJ} 309, 1994):
\begin{center}\small
\begin{tabular}{@{}lrr@{}}
\toprule
 & Treatment A & Treatment B \\
\midrule
Small stones & $93\%\ (81/87)$ & $87\%\ (234/270)$ \\
Large stones & $73\%\ (192/263)$ & $69\%\ (55/80)$ \\
Overall & $78\%\ (273/350)$ & $83\%\ (289/350)$ \\
\bottomrule
\end{tabular}\end{center}
Explain why the comparison between the treatments reverses when the groups are combined.
Can these data alone establish which treatment is better?
\end{probleminteresting}
\begin{probleminteresting}[Reservoir sampling]
Let $A$ be an array of $n\ge3$ distinct elements, indexed starting at 1.
The following algorithm is executed, with independent random choices in different iterations:
\begin{algorithm}[H]
\caption{Reservoir sampling for one element}
$sample \gets A[1]$\;
\For{$i \gets 2$ \KwTo $n$}{
  $j \gets$ a random integer, uniformly distributed on $\{1,\dots,i\}$\;
  \If{$j = 1$}{ $sample \gets A[i]$\; }
}
\Return{$sample$}
\end{algorithm}
What is the probability that the final value of $sample$ is $A[1]$? $A[3]$? $A[n]$?
Prove that the algorithm selects an element of the array uniformly.
\end{probleminteresting}
\begin{problem}
There are $n\ge1$ urns, each containing $m$ white and $k$ black balls, where $m+k>0$.
We transfer a uniformly chosen ball from the first urn to the second; after adding it,
we choose a ball uniformly from the current contents of the second urn and transfer
it to the third, and so on. Find the probability that a ball drawn uniformly from
the final urn at the end is white.
\end{problem}
\subsection*{Challenges and independent research}
\begin{problemstar}[A hidden common cause]
Let $E$ be an event with $\Pbb(E)=q$, where $0<q<1$, and let $A$ and $B$ be
conditionally independent given the value of $\ind_E$. Let
\[
\Pbb(A\mid \ind_E=e)=a_e
\quad\text{and}\quad
\Pbb(B\mid \ind_E=e)=b_e,
\qquad\text{for }e\in\{0,1\}.
\]
\begin{enumeratea}
  \item Prove that
  \[ \Cov(\ind_A,\ind_B)=q(1-q)(a_1-a_0)(b_1-b_0). \]
  When is this covariance zero? Explain why, for indicators, this is equivalent
  to independence of $A$ and $B$.
  \item Assuming $\Pbb(A\cap B)>0$, find $\Pbb(E\mid A\cap B)$.
  \item Let
  \[ q=0.10,\qquad a_0=0.02,\quad a_1=0.20,\qquad b_0=0.03,\quad b_1=0.25. \]
  Calculate $\Pbb(A)$, $\Pbb(B)$ and $\Cov(\ind_A,\ind_B)$. Find and compare
  \[ \Pbb(E),\qquad\Pbb(E\mid A),\qquad\Pbb(E\mid B),\qquad\Pbb(E\mid A\cap B). \]
  Check that
  \[ \Pbb(A\cap B)=0.00554\qquad\text{and}\qquad\Pbb(E\mid A\cap B)=\frac{250}{277}. \]
  Interpret the results by viewing $E$ as a hidden common cause and $A,B$ as two observed clues to it.
\end{enumeratea}
\end{problemstar}
\begin{problemstar}[An urn of unknown composition]
For a fixed $n\ge2$, choose $K$ uniformly from $\{0,1,\dots,n\}$. Then use an
urn containing $K$ red and $n-K$ blue balls. Make two draws from it.
Consider separately sampling with and without replacement. Let $R_i$ be the event
``the $i$th ball drawn is red'', for $i\in\{1,2\}$.
\begin{enumeratea}
  \item Given that the first ball is red, prove that in both models
  \[ \Pbb(K=k\mid R_1)=\frac{2k}{n(n+1)},\qquad\text{for }k\in\{0,\dots,n\}. \]
  \item Find $\Pbb(R_2\mid R_1)$ in both models and prove that the results are, respectively,
  \[ \frac{2n+1}{3n}\qquad\text{and}\qquad\frac23. \]
  \item Explain probabilistically why the results differ for finite $n$ and why
  the difference disappears as $n\to\infty$.
\end{enumeratea}
\end{problemstar}
\begin{researchtask}[When does a repeated test fail to provide independent evidence?]
Let $D\sim\Ber(0.01)$ indicate the presence of a disease, and let $B\sim\Ber(q)$
indicate a problematic laboratory batch, with $D$ and $B$ independent. All tests
on a given person use the same batch. Conditional on $(D,B)$, the results
$T_1,T_2,\dots$ are independent and identically distributed, with $T_i=1$ meaning a positive result.
Initially let $q=0.10$, with positive-result probabilities given by
\[
\begin{array}{c|cc}
 & B=0 & B=1\\ \hline
\Pbb(T_i=1\mid D=1,B) & 0.95 & 0.80\\
\Pbb(T_i=1\mid D=0,B) & 0.05 & 0.30
\end{array}
\]
\begin{enumeratea}
  \item Derive an exact formula for
  \[ p_k=\Pbb(D=1\mid T_1=\cdots=T_k=1),\qquad\text{for }k\ge1. \]
  \item Compare it with the naive model that first averages sensitivity and the
  false-positive probability over $B$, then assumes that test results are
  conditionally independent given $D$. Explain why this assumption can change the answer.
  \item Make a table for $k=1,\dots,5$. Check that
  \[ p_1=\frac{17}{152},\qquad p_2=\frac{701}{1592}\approx0.4403, \]
  whereas the naive model gives approximately $0.6109$ for $k=2$.
  \item Consider the parameter grid
  \[
  q\in\{0,0.05,0.10,0.20\},\qquad
  \Pbb(T_i=1\mid D=0,B=1)\in\{0.10,0.20,0.30,0.40\},
  \]
  keeping all other parameters fixed. Find the pairs for which the naive estimate
  of $p_2$ exceeds the exact value by at least $0.10$.
  \item Does the naive model always overestimate $p_2$ if all four probabilities
  in the table may vary? Give a proof or a counterexample and explain the result.
  \item Does the naive model always overestimate $p_2$ if all four probabilities
  in the table may vary? Find a necessary and sufficient condition for overestimation,
  give a counterexample, and interpret the condition.
\end{enumeratea}
\end{researchtask}
\noindent\textbf{Lean.} Let $A,B\subseteq\Omega$, with indicators taking values
in $\mathbb N$. Use the definition from the first topic and prove in Lean 4 with
Mathlib that, for every $x\in\Omega$,
\[ \ind_{A\cap B}(x)=\ind_A(x)\ind_B(x). \]
